\section{Background and Threat Model}

\subsection{Agent Runtime Model}

PRISM is designed for a tool-augmented gateway runtime in which an LLM-driven agent interacts with users, constructs prompts, invokes tools, processes tool-returned content, and emits externally visible responses. In this setting, security-relevant information does not flow through a single input channel. Instead, it traverses several stages that are all potentially attackable: user messages, retrieved external content, intermediate prompt construction, tool invocation parameters, tool-returned artifacts, outbound messages, and selected local control files.

The runtime entities relevant to PRISM are the user, the OpenClaw gateway, the in-process PRISM plugin, and a set of optional sidecar services. The sidecars include a scanner service for remote classification of suspicious tool-returned content, an invoke-guard proxy for policy-backed tool governance, a dashboard for audit inspection and policy management, and a file monitor for integrity checks over selected local files. The plugin serves as the point of lifecycle interception inside the gateway; the sidecars extend scanning, enforcement, operations, and monitoring capabilities without requiring a fork of the upstream framework.

At a high level, the runtime proceeds as follows. A user message enters the gateway and may be inspected before prompt construction. The gateway then prepares the model-facing prompt, during which security context may be injected if earlier stages have accumulated sufficient risk. The agent may subsequently invoke tools, whose parameters and destinations can be checked before execution. Tool results may then be scanned, sanitized, persisted, or incorporated into later model context. When the gateway prepares an outbound response, the final content can be checked again for residual injection signals, secret leakage, or elevated conversation risk before it is sent externally. This staged execution model is the background assumption that motivates PRISM's lifecycle-wide defense design.

\subsection{Protected Assets}

We identify five classes of assets as security-relevant within this runtime model. The first is hidden control context, including system prompts and other internal instructions that shape agent behavior but are not intended to be exposed or overridden by untrusted content. The second is tool execution privilege: the ability to invoke tools, reach selected network destinations, or perform actions whose effect extends beyond text generation. The third is local configuration and control state, including security-relevant files whose modification could weaken the gateway's policy boundary. The fourth is secret material, such as API tokens, credentials, and other sensitive values that could be disclosed through tool results or outbound messages. The fifth is security telemetry itself, including audit records and mutable policy state that operators rely on to inspect and adjust the system.

These assets matter because attacks against tool-using agents are rarely limited to prompt misclassification. A successful adversary may instead attempt to induce unsafe tool execution, cause the agent to disclose sensitive material, poison persistent state, or tamper with the control and audit surfaces that operators depend on during incident response and recovery.

\subsection{Adversaries and Goals}

PRISM considers adversaries with four principal capabilities. First, an attacker may directly control user-provided text and attempt to override instructions, exfiltrate hidden context, or induce unsafe tool use through ordinary conversational input. Second, an attacker may control external content consumed by the agent---such as fetched web pages, API responses, or other tool-returned text---and use that content to deliver indirect prompt injection after the original user message has already passed inspection. Third, an attacker may attempt to steer outbound behavior by causing the agent to transmit sensitive content to risky or untrusted destinations. Fourth, a local attacker with file-level access to selected paths may attempt to modify protected control files or interfere with the evidence trail used for auditing and incident response.

Within this model, the adversary goals that PRISM addresses are prompt override, policy bypass, unsafe tool execution, secret exfiltration, long-horizon escalation through repeated low-grade signals, and control-file tampering. These attack classes overlap with practitioner-oriented taxonomies such as the OWASP Top 10 for LLM Applications and MITRE ATLAS, but PRISM narrows them to the subset that a gateway-attached runtime defense can realistically mediate \cite{owaspllmtop102025,mitreatlas2023}. The long-horizon case is particularly important because some attacks do not appear decisively malicious at any single checkpoint. Instead, they emerge as a pattern of suspicious signals distributed across multiple turns, tools, or external interactions. This observation is one reason the paper emphasizes risk accumulation and thresholded runtime response rather than single-shot classification alone.

\subsection{Out-of-Scope Attacks and Assumptions}

PRISM does not attempt to address every security problem relevant to LLM agents. We treat model poisoning, training-time corruption, and arbitrary compromise of model internals as out of scope. We also exclude full host compromise, kernel- or hardware-level attacks, and arbitrary supply-chain compromise outside the PRISM deployment boundary. PRISM is not intended to serve as a complete OS sandbox or a general-purpose outbound firewall for all processes on the host.

The paper therefore assumes that the gateway process, the host environment, and the deployment stack retain a baseline level of integrity beyond what PRISM itself enforces. Under that assumption, PRISM aims to improve runtime security at the agent-gateway layer by making attacks more detectable, more interruptible, and more auditable across the stages it directly mediates.
